% ==================================================
% PREAMBLE
% ==================================================
\documentclass[12pt]{article}

% -------- Packages --------
\usepackage[utf8]{inputenc}
\usepackage{natbib}
\usepackage{setspace}
\usepackage[english]{babel}
\usepackage{amsmath}
\usepackage{graphicx}
\usepackage[colorlinks=true, allcolors=blue]{hyperref}
\hypersetup{
    colorlinks=true,
    linkcolor=black,
    filecolor=magenta,      
    urlcolor=cyan,
    pdfpagemode=FullScreen
}
\usepackage{wasysym}
\usepackage{mathabx}
\usepackage{enumitem}
\usepackage{comment}
\usepackage[table]{xcolor}
\usepackage{geometry}
\geometry{legalpaper,margin=0.7in}    % margin left and right is 1inch
% \usepackage[legalpaper,top=1cm,bottom=2cm,left=0.5cm,right=0.5cm,marginparwidth=1.75cm]{geometry}
\usepackage{wrapfig}
\usepackage{caption}
\usepackage{subcaption}
\usepackage{multicol}
\usepackage{fancyhdr}
\usepackage{datetime2}   % better date formatting
\usepackage{tabularx}
\usepackage{multirow}
\usepackage{array}
\usepackage{makecell}
\usepackage[most]{tcolorbox}
\usepackage{arydshln}

% -------- Custome Commands --------

\pagestyle{fancy}
\fancyhf{}  % clear default header/footer
\fancyhead[L]{COHA Cheatsheet} % header content
\fancyhead[R]{Last edited: \today} % header content
\fancyfoot[R]{\thepage} % page number in footer

\DTMsetdatestyle{ddmmyyyy}

\newcounter{row}
\newcolumntype{N}{>{\stepcounter{row}\arabic{row}.}c}

\setlength{\parindent}{0pt}

\renewcommand{\arraystretch}{1.5}

%%%%%%%%%%%%%%%%%%%%%%%%%%%%%%%%%%%%%%%%%%%%%%%%%%%%
% HEADER + TOC
%%%%%%%%%%%%%%%%%%%%%%%%%%%%%%%%%%%%%%%%%%%%%%%%%%%%

\begin{document}
{
  \hypersetup{linkcolor=black}
  \tableofcontents
}
\clearpage

%%%%%%%%%%%%%%%%%%%%%%%%%%%%%%%%%%%%%%%%%%%%%%%%%%%%
% CONTENT
%%%%%%%%%%%%%%%%%%%%%%%%%%%%%%%%%%%%%%%%%%%%%%%%%%%%
\subsection{Generalinfo}
\begin{itemize}
    \item Telefonnummern von Leitstand und Parkkassa $\longrightarrow$ siehe Signal-Gruppenbeschreibung
    \item Gerätelogins und Schlüsselcodes $\longrightarrow$ siehe Signal-Gruppenbeschreibung
    \item Ready2Order-App (R2O) \textbf{muss am Tablet immer offen} sein, damit die Bestellungen beim Drucker hinten raus kommen.
\end{itemize}

\subsection{ToDos und Diskussionspunkte}
\begin{itemize}
    \item brauch ma Anleitung zu Schank/Kaffeemaschine/Milchsystem reinigen?
    \item Anleitung zu Abrechnung machen?
    \item Ab wann darf Kaffeemaschine/Schank geschlossen werden? 17:45 Uhr in Ordnung? (Wenn viel los ist und noch bis 18 Uhr verkauft werden kann ist klar dass man noch bis 18 Uhr verkauft. Aber wenn wenig los und man potenziell nur ein-zwei Personen wegschickt weil Kaffeemaschine schon geschlossen ist, wäre ok?)
    \item Darf Gastgarten auch erst ab 18 Uhr geschlossen werden? (Wieder außer es ist nichts los)
    \begin{itemize}
        \item[] $\longrightarrow$ wenn nur wenig vor 18 Uhr erledigt werden darf, dann dauert der Abenddienst sehr lange
    \end{itemize}
    \item TODO: Service station im Gastgarten?
    \item TODO: unnützes Zeug aus Lager räumen: Kisten mit catering zeug, Eis-zubehör (Waffeln, Servietten)
    \item TODO: Produktlisten ausdrucken und einfolieren um Bestands-check zu vereinfachen?
    \item Rezept für CreamCheese Bagel?
\end{itemize}


\clearpage
%%%%%%%%%%%%%%%%%%%%%%%%%%%%%%%%%%%%%%%%%%%%%%%%%%%%
\section{Produktlisten}

\begin{tabularx}{\textwidth}{|c|X r|c|}
\hline
\rowcolor{gray!25}
\multicolumn{2}{|l}{\large{\textbf{Getränke - Alkoholfrei}}\rule{0pt}{2.5ex}} & Info & --- \\ \hline
%
\multirow{5}{*}{\rotatebox[origin=c]{90}{Fritz Kola}}
& Rhabarber & = FK Saison & \\
& Orange && \\
& Mischmasch & wie spezi& \\
& Classic && \\
& Light && \\ \hline
%
\multirow{5}{*}{\rotatebox[origin=c]{90}{Bauernsäfte}}
& Apfel & von hasenfit& \\
& Orange & hasenfit& \\
& Marille & hasenfit& \\
& Johannisbeere & hasenfit& \\
& Pfirsich & nur für hausgemachten Eistee& \\ \hline
%
\multirow{4}{*}{\rotatebox[origin=c]{90}{Limos}}
& Makava & Zitrone& \\
& Frucade && \\
& Tonic & von RedBull& \\
& Rote Rakete & koffeinhaltig, von Paul\&Bohne& \\ \hline
%
\multirow{3}{*}{\rotatebox[origin=c]{90}{Sirup}}
& Holunder & von Biohof Greiml am Schoberpass& \\
& Zitronensaft & von Rauch& \\
&&& \\ \hline
%
\multirow{3}{*}{\rotatebox[origin=c]{90}{Wasser}}
& prickelnd 0.33 & von vöslauer& \\
& still 0.33 & von vöslauer& \\
&& & \\ \hline
\end{tabularx}

\vspace{1em}

\begin{tabularx}{\textwidth}{|c|X r|c|}
\hline
\rowcolor{gray!25}
\multicolumn{2}{|l}{\large{\textbf{Getränke - Alkohol}}\rule{0pt}{2.5ex}} & Info & ---\\ \hline
%
\multirow{3}{*}{\rotatebox[origin=c]{90}{Bier}}
& Gösser Biostoff & Fass & \\
& Bier af & ZeroLabs/Gösser Naturgold 0.0\%/Moretti & \\
& Craftbeer & von Tom\&Harrys & \\ \hline
%
\multirow{6}{*}{\rotatebox[origin=c]{90}{Wein}}
& Sauvignon Blanc & & \\
& Gelber Muskateller & & \\
& Weißburgunder & & \\
& Prosecco & im Fass & \\
& Sekt & Flasche, Brut, Harkamp & \\
& Rot & & \\ \hline
%
\multirow{5}{*}{\rotatebox[origin=c]{90}{Spirituosen}}
& Aperol &&\\
& Gin &&\\
& Vodka &&\\
& Rum &&\\
&&&\\ \hline
\end{tabularx}

\vspace{1em}

\begin{tabularx}{\textwidth}{|c|X r|c|}
\hline
\rowcolor{gray!25}
\multicolumn{2}{|l}{\large{\textbf{Getränke - Heiß/Iced}}\rule{0pt}{2.5ex}} & Info & --- \\ \hline
%
\multirow{4}{*}{\rotatebox[origin=c]{90}{Milch}}
& Kuhmilch & in Eimer &\\
& Hafermilch & von Oatly&\\
& Laktosefrei &&\\
& Haltbarmilch & nur als Notfall wenn andere aus&\\ \hline
%
\multirow{3}{*}{\rotatebox[origin=c]{90}{Kaffee}}
& Kaffeebohnen & Paul\&Bohne &\\
& Decaf & Paul\&Bohne &\\
& & &\\ \hline
%
\multirow{11}{*}{\rotatebox[origin=c]{90}{Tee}}
& Kräuter &&\\
& Schwarztee - Chai &&\\
& Schwarztee - Assam &&\\
& Schwarztee - Earl Grey &&\\
& Schwarztee & der von Hornig am ehesten verwenden für Eisteeansatz &\\
& Grüntee &&\\
& Früchte &&\\
& Kamille &&\\
& Lindenblüten &&\\
& Pfefferminze &&\\
& Chai &&\\ \hline
%
\multirow{2}{*}{\rotatebox[origin=c]{90}{}}
& Kakao &&\\
& Schlag & gut schütteln vor einfüllen&\\ \hline
\end{tabularx}

\vspace{1em}

\begin{tabularx}{\textwidth}{|c|X r|c|}
    \hline
    \rowcolor{gray!25}
    \multicolumn{2}{|l}{\large{\textbf{Verbrauchsmaterial}}\rule{0pt}{2.5ex}} & Info & ---\\
    \hline
    %
    \multirow{3}{*}{\rotatebox[origin=c]{90}{}}
    & Servietten &&\\
    & Strohhalme &&\\
    & Zahnstocher &&\\ \hline
    %
    \multirow{4}{*}{\rotatebox[origin=c]{90}{ToGo}}
    & Einwegbecher &&\\
    & Mehrwegbecher &&\\
    & Becher-Deckel &&\\
    & Papierteller &&\\ \hline
    %
    \multirow{4}{*}{\rotatebox[origin=c]{90}{}}
    & Gaspatronen &&\\
    & Abrechnungszettel &&\\
    & Stundenlisten &&\\
    & Coha-Papier &&\\ 
    & Bon-rollen && \\ \hline
    %
    \multirow{6}{*}{\rotatebox[origin=c]{90}{Reinigung}}
    & frische Geschirrhangerl und Fetzen &&\\ 
    & Müllsäcke &&\\
    & Desinfektionsmittel & &\\
    & Glasreiniger & &\\ 
    & Geschirrspülmittel & &\\ 
    & Kaffee-Reinigungspulver & &\\ \hline
\end{tabularx}

\begin{tabularx}{\textwidth}{|c|X r|c|}
\hline
\rowcolor{gray!25}
\multicolumn{2}{|l}{\large{\textbf{Essen}}\rule{0pt}{2.5ex}} & Info & ---\\ \hline
%
\multirow{2}{*}{\rotatebox[origin=c]{90}{R\&F}}
& Bagels &&\\
& Butter Croissants &&\\ \hline
%
\multirow{11}{*}{\rotatebox[origin=c]{90}{Bagel-Zutaten}}
& Frischkäse &&\\
& Pesto &&\\
& Schinken &&\\
& Käse &&\\
& Tomaten &&\\
& Gurke &&\\
& Mozzarella &&\\
& Lachs &&\\
& Senf &&\\
& Rucola &&\\
& Zupfsalat &&\\ \hline
%
\multirow{3}{*}{\rotatebox[origin=c]{90}{Deko}}
& Minze & für Eistee, Hugo und evtl. Holundersoda &\\
& Orangen &&\\
& Zitronen &&\\ \hline
%
\multirow{4}{*}{\rotatebox[origin=c]{90}{Eiscreme}}
& Vanille & nur für Eiskaffe und Affogato &\\
& Twinnie &&\\
& Jolly &&\\
& Magnum &&\\ \hline
%
\multirow{4}{*}{\rotatebox[origin=c]{90}{Sonstiges}}
& Zucker &&\\
& Süssstoff &&\\
& Ingwer & für hausgemachten Eistee &\\ 
&&& \\ \hline
\end{tabularx}


\clearpage
%%%%%%%%%%%%%%%%%%%%%%%%%%%%%%%%%%%%%%%%%%%%%%%%%%%%
\section{Checklisten}

Reihenfolge so gewählt, dass es zeitlich Sinn macht. Zum Beispiel im Abenddienst: der Gastgarten kann zuerst zugesperrt werden und erst später die Kaffeemaschine. 
\newline \newline
Ausschank muss bis 18Uhr möglich sein. Erst danach dürfen Schank und Kaffeemaschine geputzt werden. \textcolor{red}{$\longrightarrow$ TODO}. Außer es ist ganz wenig los und absehbar dass 'Niemand' mehr kommt, dann kann schon ab 18:45Uhr angefangen werden.

\subsection{Frühdienst}
\begin{tabularx}{\textwidth}{|N|X r|}
    \hline
    & Reservierungen checken & im Kalender am iPad\\
    & Backofen, Geschirrspüler, Eisvitrine einschalten & brauchen alle Zeit bis einsatzbereit sind \\
    & Kuchenvitrine herrichten & \\
    & Bagels herrichten & zuerst die vom Vortag verwenden \\
    & Croissants aufbacken & zuerst die vom Vortag verwenden \\
    & Eisvitrine herrichten & \\
    & Gastgarten herrichten & \makecell[l]{Tische reinigen, Schirme aufspannen, \\ Tischnummern + Karten + Zuckergläser aufstellen }\\
    & gegebenenfalls Licht einschalten &\\
    & Kaffeemaschine aktivieren &\\
    & Klos aufsperren & evtl. schon von Reinigungspersonal aufgesperrt\\
    \hline
\end{tabularx}


\subsection{Abenddienst}
\setcounter{row}{0}
\begin{tabularx}{\textwidth}{|N|X r|}
    \hline
    & restliches Personalessen eintragen &\\
    & Abrechnung & erst ab 18Uhr und keine offenen Tische mehr\\
    & Gastgarten abbauen &\\
    & Getränkeladen auffüllen & \\
    & Leergut wegbringen & \\
    & Backofen ausschalten und reinigen & \\
    & Kuchenvitrine ausräumen und säubern & \\
    & Croissants und Bagels für den nächsten Tag einpacken und in Kühlladen geben. & nicht am Sonntag! \\
    & Schank putzen &\\
    & Milchsystem reinigen &\\
    & Kaffeemaschine reinigen &\\
    & Müll weg bringen, frische Müllsäcke in Eimer geben&\\
    & Oberflächen reinigen &\\
    & Eisvitrine ausschalten, Eis wieder in Lager/kleinen TK geben &\\
    & Geschirrspüler ausschalten/abpumpen/Siebe reinigen & Türe offen lassen\\
    & Licht ausschalten &\\
    & Umsatzabgabe in Briefkasten werfen &\\
    & Alle Türen + Klos zu sperren &\\
    \hline
\end{tabularx}


\subsection{Stehzeiten}
Aufgaben für unter der Woche, wenn viel Zeit ist. Wieder nach Wichtigkeit geordnet.
\begin{itemize}[itemsep=1pt, topsep=10pt]
    \item Kühlladen wischen
    \item Eisteeansatz machen
    \item Eistee machen
    \item Pflanzen gießen
    \item Kaffeemaschine Grundreinigen
    \item Glasregale wischen
    \item Eisvitrine wischen
    \item Boden kehren
\end{itemize}


\clearpage
%%%%%%%%%%%%%%%%%%%%%%%%%%%%%%%%%%%%%%%%%%%%%%%%%%%%
\section{Anleitungen}

\subsection{Bagels}

\begin{tabularx}{\textwidth}{|l|X|}
    \hline
    \rowcolor{gray!25}
    \multicolumn{2}{|l||}{\large{Schinken-Käse}\rule{0pt}{2.5ex}}\\
    \hline
    Frischkäse & auf Deckel und Boden \\
    Zupfsalat & \\
    Schinken & 2 Blätter \\
    Käse & 1 Scheibe \\
    Gurke & ca. 3 Scheiben\\
    Tomaten & ca. 3 Scheiben \\ 
    Salz \& Pfeffer & \\ \hline
\end{tabularx}

\vspace{1em}

\begin{tabularx}{\textwidth}{|l|X|}
    \hline
    \rowcolor{gray!25}
    \multicolumn{2}{|l|}{\large{Käse-Gurke}\rule{0pt}{2.5ex}} \\
    \hline
    Frischkäse & auf Deckel und Boden \\
    Zupfsalat & \\
    Gurke & 5-6 Scheiben \\
    Käse & 2 Scheiben \\
    Salz \& Pfeffer & \\ \hline
\end{tabularx}

\vspace{1em}

\begin{tabularx}{\textwidth}{|l|X|}
    \hline
    \rowcolor{gray!25}
    \multicolumn{2}{|l|}{\large{Lachs}\rule{0pt}{2.5ex}} \\
    \hline
    Frischkäse & auf Deckel und Boden \\
    Senf & Teelöffel auf Deckel od. Boden verstreichen \\
    Gurke & ca. 5 Scheiben \\
    Lachs & ca. 2 Stk. \\
    Rucola & \\
    Salz \& Pfeffer & \\ \hline
\end{tabularx}

\vspace{1em}

\begin{tabularx}{\textwidth}{|l|X|}
    \hline
    \rowcolor{gray!25}
    \multicolumn{2}{|l|}{\large{Tomaten - Mozzarella}\rule{0pt}{2.5ex}} \\
    \hline
    Pesto & auf Deckel und Boden \\
    Tomaten & 4-5 Scheiben \\
    Mozzarella & 2 Scheiben \\
    Rucola & \\
    Salz \& Pfeffer & \\ \hline
\end{tabularx}

\vspace{1em}

\begin{tabularx}{\textwidth}{|l|X|}
    \hline
    \rowcolor{gray!25}
    \multicolumn{2}{|l|}{\large{Creamcheese}\rule{0pt}{2.5ex}} \\
    \hline
    Frischkäse & auf Deckel und Boden \\
    Käse & 2 Scheiben \\
    Rucola & \\
    Salz \& Pfeffer & \\ \hline
\end{tabularx}

\clearpage
\subsection{Kaffee \& co.}

\begin{tabularx}{\textwidth}{|X|c|c|c|c|l|X|}
    \hline
    Typ &
    Espresso &
    \makecell{Wasser\\(y/n)} &
    \makecell{Milch \\ (y/n)} &
    \makecell{Schaum \\ (y/n)} &
    Tasse &
    Zusatz-infos\\
    \hline
    Espresso    & 1     & n & n & n & klein     & ='Kurzer' \\ \hdashline
    Doppio      & 2     & n & n & n & mittel    & \\ \hdashline
    Macchiato   & 1     & n & n & y & klein & \\ \hdashline
    Verlängerter & 1    & y & n & n & mittel    & \\ \hdashline
    Milchkaffe  & 1     & y & y & n & mittel    & \\ \hdashline
    Cappuccino  & 1     & n & y & y & groß      & \\ \hdashline
    Melange     & 1     & y & y & y & groß      & \\ \hdashline
    Flat White  & 2     & n & y & y & groß      & =cappu+extra shot \\ \hdashline
    Cafe Latte  & 1     & n & y & y/n & Glas mittel & mit Schaum nicht so genau \\ \hdashline
    großer Brauner & 2  & n & y & n & mittel    & normal Milch extra \\ \hdashline
    kleiner Brauner & 1 & n & y & n & klein     & normal Milch extra \\ \hdashline
    Eiskaffe    & 1     & y & n & n & Glas groß & +2Kgl Vanille Eis +Schlag \\ \hdashline
    Affogato    & 1     & n & n & n & mittel    & +1Kgl Vanille Eis \\ \hdashline
    Kakao       & 0     & y & y & y/n & Glas mittel & 1,5TL +Schlag \\ \hdashline
    Chai Latte  & 0     & y & y & y/n & Glas mittel & beutel 3min ziehen \\
    \hline
\end{tabularx}

\vspace{1em}

Some Extra Info:
\begin{description}[itemsep=1pt, topsep=0pt]
    \item[Mocca:] $\neq$ Espresso, wird mit feiner gemahlenen Kaffee-Pulver am Herd in Kupferkanne gekocht 
    \item[Dirty Chai/Kakao =] Chai/Kakao + Espresso shot
    \item[Ristretto]
    \item[Cortado]
    \item[Espresso Lungo]
    \item[Wiener Melange]
\end{description}


\subsection{Eistee}
\begin{enumerate}[topsep=0pt, itemsep=1pt]
    \item Eistee - Ansatz machen
    \begin{itemize}
        \item in Paul\&Bohne Eimer
        \item bis zur Hälfte mit heißem Wasser
        \item 5-6 Schwarztee (nicht der von PureTee) hinzugeben (Papiertag wegschneiden)
        \item mind. eine Nacht im Fass-kühler ziehen lassen, kann auch länger stehen
    \end{itemize}
    \item gewaschene Milchflaschen bis zur Hälfte mit Eisteeansatz füllen
    \item Rest mit Pfirsichsaft auffüllen (aber nicht ganz voll!)
    \item in Blender 2 TL Ingwer-Konzentrat mit etwas Pfirsichsaft mixen (für ca. 3 Flaschen, wenn mehr gemacht wird mehr Ingwer nehmen) und zu den Eistees dazu geben
    \item Flaschen mit Datum beschriften
\end{enumerate}

\subsection{Smoothies}
\textcolor{red}{TODO}

\clearpage
%%%%%%%%%%%%%%%%%%%%%%%%%%%%%%%%%%%%%%%%%%%%%%%%%%%%
\section{Kuchen und Desserts}

Schokoschnitte
Karottenkuchen
Mamorkuchen
Topfenbeeren - schnitte
Zwetschgen-Polenta Törtchen
Apfel-Küchlein
Cupcakes
Nuss-Amaranth-Riegel

\section{Allergene}

\section{Reservierungen}

\section{Dienstplan und Schichtplanung}

\section{Schulungen}
\subsection{Kaffee Paul und Bohne}
\subsection{Brand-Löschübung}

\end{document}
